\documentclass[11pt,a4paper]{article}
\usepackage[T1]{fontenc}
\usepackage{lmodern,microtype}
\usepackage[a4paper,margin=25mm]{geometry}
\usepackage{amsmath,amssymb,booktabs,array}
\usepackage[dvipsnames]{xcolor}
\usepackage[colorlinks=true,urlcolor=MidnightBlue,linkcolor=MidnightBlue]{hyperref}
\setlength{\emergencystretch}{2em}
\setlength{\parskip}{3pt}
\hypersetup{pdftitle={Rank 5: publication revision letter},pdfauthor={}}
\title{Rank 5: publication revision letter}
\author{}\date{24 November 2026}
\begin{document}
\maketitle

This revision prepares the fourth reviewed manuscript for the existing
GitHub series and a subsequent Hexagon submission. The title is now
\emph{Grothendieck rings of complex fusion categories: rank 5},
matching the titles of the rank-two, rank-three and rank-four papers.
The manuscript has \textbf{72 pages and 51 bibliography entries},
including \textbf{22 additional references}. Its theorems, sharp bounds,
seventeen model tensors, certificates and verification algorithms are
unchanged. This is a publication and attribution revision, not a new
mathematical classification or a claim of formal certification.

\section*{Title, context and model conventions}

The printed title and PDF title metadata were aligned with the series.
The existing empty author field was retained. A short research-provenance
note states that AI assisted with the research and preparation and that
complete human verification remains pending.

The introduction now cites the standard textbook background and gives
a fuller, carefully qualified history. In particular, it directly credits
the pivotal rank-four classification of Dong--Palcoux--Plessis and the
pseudo-unitary premodular rank-five classification of
Bruillard--Ortiz-Marrero. It distinguishes Larson's candidate rings
from the unitary classification of Edie-Michell--Izumi--Penneys and
describes Wang--Dong's faithfully graded rank-five results as partial.
The earlier modular and premodular classifications and Hagge--Hong's
low-rank sphericality observation are also included.

The new model-source subsection specifies the quantum parameter
$q=\exp(\pi i/(k+n))$ for $\mathrm{SU}(n)_k$, the even-label adjoint
subcategory for $\mathrm{SU}(2)_k$, and the five adjoint weights at
$\mathrm{SU}(3)_4$. It assigns construction and unitarity references
to the models individually, gives the precise Liu--Palcoux--Ren
construction sections, and credits Liu and Snyder with recognizing
the $C_3+6$ example. Izumi remains the source for the concrete
$C^*$ realization and the involution used at the dimension endpoint.

\section*{The 22 added references and their roles}

\begin{center}\small
\begin{tabular}{@{}p{.37\linewidth}p{.59\linewidth}@{}}
\toprule
Source & Role in the revised manuscript\\\midrule
Etingof--Gelaki--Nikshych--Ostrik & Standard fusion-category, module-category and center background.\\
Barrett--Westbury & Original spherical-category terminology.\\
Hagge--Hong & Existence of spherical structures in pivotal categories of rank at most four.\\
Ostrik, premodular rank three & Earlier classification with a ribbon structure.\\
Rowell--Stong--Wang & Unitary modular classification through rank four.\\
Bruillard, premodular rank four & Rank-four classification up to Grothendieck equivalence.\\
Larson & Non-self-dual pseudo-unitary rank-four candidate rings.\\
Edie-Michell--Izumi--Penneys & Unitary quadratic categories and the rank-four dual-pair case.\\
Dong--Palcoux--Plessis & The pivotal fifteen-ring classification at rank four.\\
Bruillard--Ortiz-Marrero & Pseudo-unitary premodular rank-five classification.\\
Dong--Chen--Wang & Integrality of a subcategory with rank one less than the ambient rank.\\
Gelaki--Nikshych & Universal and square-class gradings.\\
M\"uger & Modularity, dimension and Gauss sums of the spherical center.\\
Ng--Schauenburg, pivotal indicators & The pivotal rotation and its finite order.\\
Ng--Schauenburg, indicators and exponents & The ordinary center trace formula and the twist-order/exponent identity.\\
Schauenburg, center and bimodules & Morita invariance of the center.\\
Tambara--Yamagami & Original construction of the two named model rings.\\
Andersen--Paradowski & Quantum-group construction and fusion rules.\\
Siehler & The affine-group near-group example at $q=5$.\\
Evans--Gannon & Explicit $C^*$ Tambara--Yamagami realizations.\\
Liu & Historical Liu--Snyder recognition of the $C_3+6$ example.\\
Burciu--Natale & Equivariantization simples and fusion formulas, using the corrected arXiv version.\\
\bottomrule
\end{tabular}
\end{center}

Each added entry is cited at a specific historical or mathematical use.
The bibliography has no duplicate keys, undefined citations or uncited
entries. The classification remains one of based rings admitting
categorification, not a classification of all tensor equivalence classes.

\section*{More precise attributions within the proofs}

Existing ENO citations now identify the pseudo-unitary spherical
structure, global-dimension comparison, sphericalization, reality of
dimensions, cyclotomicity and induction--restriction results precisely.
The formal-codegree normalization, $d$-number theorem and executable
coefficient criterion have precise Ostrik locators. The universal and
square-class gradings, modular-center normalization, Morita-center
equivalence, ordinary indicator trace formula, exponent identity and
finite-order pivotal rotation now have their original sources attached.
The short maximal-rank-subcategory argument is retained with direct
credit to Dong--Chen--Wang.

At the dimension endpoint, Izumi's existence theorem is cited separately
from the specific symmetry. The equivariantization formulas cite
Burciu--Natale, Corollary~2.13 and Theorem~3.9, in the corrected 2019
arXiv version. The broad collective realization citation at the end of
the classification proof is replaced by a reference to the individual
construction sources following Table~1. Schopieray and Meng--Yu remain
cited for the near-group restrictions they establish.

All 114 labels and all 84 theorem-like statements were compared with
the fourth delivery. Mathematical assertions and formulas are retained;
the source patch records the editorial changes and citation insertions.

\section*{GitHub links and publication materials}

The introduction links to the
\href{https://github.com/sebastienpalcoux/grothendieck-ring-classification/blob/main/manuscripts/rank5.tex}{current source},
\href{https://github.com/sebastienpalcoux/grothendieck-ring-classification/blob/main/manuscripts/rank5/README.md}{reading guide},
\href{https://sebastienpalcoux.github.io/grothendieck-ring-classification/}{project website},
and the exact
\href{https://github.com/sebastienpalcoux/grothendieck-ring-classification/tree/b479c65c702c6a09a674df7c2f8468b1a48915d3/supplements/rank5/sharp}{verification snapshot}
at commit \texttt{b479c65c702c}.
The snapshot preserves the reviewed programs, certificates, tensors and
actual verification transcripts. The current canonical manuscript is
\path{manuscripts/rank5.tex}, with a matching PDF; the supplement contains
matching reading copies and the earlier review history.

The repository homepage and project website present ranks two through
five uniformly. Rank five is now described by the unrestricted theorem,
with seventeen based rings, sphericality, a pseudo-unitary Galois
conjugate, and the sharp bounds $\mu\le4$ and
$\operatorname{FPdim}\le48+24\sqrt3$. The older multiplicity-twenty
census, its manuscript and its verification guide retain their bytes
and are identified as historical supporting material. The higher-rank
decision records and six unresolved cases are unchanged.

The publication includes a one-file Hexagon source archive containing
\path{main.tex}, a complete portable source-and-verification archive,
this letter, an exact manuscript patch and refreshed file identities.
The three earlier Hexagon source archives are retained unchanged.
The rank-five files are prepared for submission; this revision does not
claim that a Hexagon deposit has already been made or accepted.

\end{document}
